\section{2026-08-22 — Dos entrevistas, y un módulo que sobrevive por la razón equivocada}

Con las dos entrevistas realizadas quedó cerrada la sección de user research y desaparecieron los últimos \enquote{Completar.} del capítulo de descripción. Conviene registrar tres cosas: una que salió mejor de lo esperado, una que obligó a decidir, y una que quedó floja y va a ser preguntada.

Lo que salió mejor de lo esperado es la triangulación. Tres hallazgos aparecen de forma independiente en la encuesta y en las dos entrevistas, sin que ninguna de las preguntas los indujera. El primero es que la evidencia no es una prestación sino una condición de permanencia: la periodista lo formuló en términos de abandono, dijo que dejaría de usar la herramienta si no pudiera acceder a la fuente empleada. El segundo es que el riesgo que se percibe no es el error técnico sino la toma de partido. El tercero es que el producto se entiende como asistente y no como oráculo. Ninguno de los tres instrumentos por separado sostenía esas afirmaciones con la fuerza con que las sostienen los tres juntos, y conviene apoyar la defensa en eso antes que en cualquier métrica del clasificador.

Lo que obligó a decidir fue una contradicción entre las dos entrevistas sobre el módulo de credibilidad de la cuenta autora. El fundador, consultado sobre qué recortaría ante una restricción de plazo, señaló ese módulo con un argumento sólido: es el componente donde resulta más fácil introducir sesgo y más difícil justificar la decisión, porque atribuir credibilidad por cantidad de seguidores o por distintivo de verificación es arbitrario y, sobre todo, manipulable, dado que ambas cosas se compran. La periodista, en cambio, describió la observación de la cuenta como el primer paso de su proceso diario, y enumeró exactamente las señales que el otro entrevistado cuestiona.

La resolución consistió en separar dos dimensiones que hasta ahora se venían tratando como una sola: el valor de uso y la validez probatoria. Las señales de la cuenta tienen valor de uso alto, porque el profesional las consulta primero y el setenta y uno por ciento del segmento primario también, pero tienen validez probatoria baja. En consecuencia el módulo se conserva como señal exhibida al usuario y se le reduce el peso en la determinación del puntaje, que pasa a apoyarse en el clasificador y en el contraste con evidencia externa. Vale dejar asentado que el módulo sobrevive por su valor informativo y no por su capacidad probatoria, porque es una distinción que no estaba hecha en ninguna parte del trabajo y que ordena varias decisiones posteriores.

La entrevista al fundador dejó además una precisión sobre el argumento competitivo que conviene incorporar a la defensa. El capítulo de descripción sostenía que la ventaja no reside en la técnica sino en el modelo de negocio. Consultado sobre por dónde atacaría el proyecto si compitiera con él, el entrevistado descartó la vía tecnológica y respondió que competiría por distribución y por confianza percibida. Esos dos componentes estaban implícitos en la formulación original y ahora quedan explicitados.

Lo que quedó flojo es el perfil de la segunda entrevista. El diseño original preveía un fact-checker o un académico de desinformación, y lo que se consiguió fue un fundador de una empresa de inteligencia artificial. El reemplazo aporta autoridad sobre la tecnología y sobre el negocio, que es donde el trabajo estaba más expuesto, pero no aporta nada sobre la metodología de verificación ni sobre el contexto mediático argentino. Ninguna de las dos entrevistas cubre a una organización de verificación profesional. Es la pregunta previsible del tribunal y conviene llegar con la respuesta preparada, que es que la gestión se inició y quedó fuera de plazo.

Quedan dos puntos abiertos que también conviene tener presentes. La periodista advirtió sobre las afirmaciones que combinan datos verdaderos con interpretación u omisión, que no admiten clasificación binaria. El esquema pasó de cuatro clases a tres eliminando precisamente la categoría intermedia, según se registró en la entrada anterior, y la clase de ausencia de verificación absorbe esos casos solo de forma parcial. Y la pregunta sobre esa misma clase, formulada al fundador, no obtuvo respuesta durante la entrevista, de modo que sigue sin validación externa.

\bigskip
